% concatenated from manuscript.tex
%% 
%% Copyright 2019-2024 Elsevier Ltd
%% 
%% This file is part of the 'CAS Bundle'.
%% --------------------------------------
%% 
%% It may be distributed under the conditions of the LaTeX Project Public
%% License, either version 1.3c of this license or (at your option) any
%% later version.  The latest version of this license is in
%%    http://www.latex-project.org/lppl.txt
%% and version 1.3c or later is part of all distributions of LaTeX
%% version 1999/12/01 or later.
%% 
%% The list of all files belonging to the 'CAS Bundle' is
%% given in the file `manifest.txt'.
%% 
%% Template article for cas-sc documentclass for 
%% double column output.

\documentclass[a4paper,fleqn]{cas-sc}

% If the frontmatter runs over more than one page
% use the longmktitle option.

%\documentclass[a4paper,fleqn,longmktitle]{cas-sc}

\usepackage[numbers]{natbib}
%\usepackage[authoryear]{natbib}
%\usepackage[authoryear,longnamesfirst]{natbib}
\usepackage{cleveref}
\usepackage{amsthm}  % already loaded by cas-sc, but safe to include
\usepackage{algorithm}%%%%%%%%%%%%%
%\usepackage{algorithmic}%%%%%%%%
\usepackage{algpseudocode}%%%%%%%%%%%%
\usepackage{float}
\newtheorem{theorem}{Theorem}[section]      % Theorem numbering per section
\newtheorem{lemma}[theorem]{Lemma}          % Lemmas share the same counter
\newtheorem{proposition}[theorem]{Proposition}
\newtheorem{corollary}[theorem]{Corollary}

\theoremstyle{definition}
\newtheorem{definition}[theorem]{Definition}

\theoremstyle{remark}
\newtheorem{remark}[theorem]{Remark}

%%%Author macros
\def\tsc#1{\csdef{#1}{\textsc{\lowercase{#1}}\xspace}}
\tsc{WGM}
\tsc{QE}
%%%

% Uncomment and use as if needed
%\newtheorem{theorem}{Theorem}
%\newtheorem{lemma}[theorem]{Lemma}
%\newdefinition{rmk}{Remark}
%\newproof{pf}{Proof}
%\newproof{pot}{Proof of Theorem \ref{thm}}
%%%%%%%%%%%%%%%%%%%%%%%%%%%%%%Run this if you are facing environment issue%%%%%%%%%%%%%%%%%
%%%%%%%%%%%%%%%%%%%%%%%%%%%%%%%%%%%%%%%%%%%%%%%%%%%%%%%%%%%%%%%%%%%%%%%%%%%%%%%%%%%%%%%%%%%%%%%%%%%
% \usepackage{biblatex}
% \addbibresource{references.bib}
%%%%%%%%%%%%%%%%%%%%%%%%%%%%%%%%%%%%%%%%%%%%%%%%%%%%%%%%%%%%%%%%%%%%%%%%%%%%%%%%%%%%%%%%%%%%%%%%%%%%
\begin{document}
\let\WriteBookmarks\relax
\def\floatpagepagefraction{1}
\def\textpagefraction{.001}

% Short title
\shorttitle{Direct Solver for Scaled Partial-Isometric Linear Systems}

% Short author
\shortauthors{S. Kar, M. Venkatapathi}  

% Main title of the paper
\title[mode=title]{An Optimal Least-Square Solver For Scaled Partial-Isometric Linear Systems}  

% Title footnote mark
% eg: \tnotemark[1]
\tnotemark[1] 

% Title footnote 1.
% eg: \tnotetext[1]{Title footnote text}
\tnotetext[1]{This work was funded by ......} 

% First author
%
% Options: Use if required
% eg: \author[1,3]{Author Name}[type=editor,
%       style=chinese,
%       auid=000,
%       bioid=1,
%       prefix=Sir,
%       orcid=0000-0000-0000-0000,
%       facebook=<facebook id>,
%       twitter=<twitter id>,
%       linkedin=<linkedin id>,
%       gplus=<gplus id>]

\author[1]{Suvendu Kar}%[<options>]

% Corresponding author indication
%\cormark[1]

% Footnote of the first author
\fnmark[1]

% Email id of the first author
\ead{sk73191491@gmail.com}

% URL of the first author
\ead[url]{}

% Credit authorship
% eg: \credit{Conceptualization of this study, Methodology, Software}
\credit{}

% Address/affiliation
\affiliation[1]{organization={Department of Computational and Data Sciences},
            addressline={Indian Institute of Science }, 
            city={Bangalore},
%          citysep={}, % Uncomment if no comma needed between city and postcode
            postcode={560012}, 
            state={karnataka},
            country={India}}

\author[1]{Murugesan Venkatapathi}%[]

\cormark[1]
% Footnote of the second author
\fnmark[2]

% Email id of the second author
\ead{murugesh@iisc.ac.in}

% URL of the second author
\ead[url]{}

% Credit authorship
\credit{}

% % Address/affiliation
% \affiliation[1]{organization={Department of Computational and Data Sciences},
%             addressline={Indian Institute of Science}, 
%             city={Bangalore},
% %          citysep={}, % Uncomment if no comma needed between city and postcode
%             postcode={560012}, 
%             state={Karnataka},
%             country={India}}

% Corresponding author text
\cortext[1]{Corresponding author}

% Footnote text
\fntext[1]{}

% For a title note without a number/mark
%\nonumnote{}

% Here goes the abstract
\begin{abstract}
We present an $O(mn)$ direct least-squares solver for $m \times n$ linear systems with a scaled partial isometry. The proposed algorithm is also useful when the system is block diagonal and each block is a scaled partial isometry with distinct scaling factors. We also include numerical experiments as a demonstration.
\end{abstract}

% Use if graphical abstract is present
%\begin{graphicalabstract}
%\includegraphics{}
%\end{graphicalabstract}

% Research highlights
%\begin{highlights}
%\item 
%\item 
%\item 
%\end{highlights}


% Keywords
% Each keyword is seperated by \sep
\begin{keywords}
 \sep Direct solver for linear systems \sep Scaled partial isometries  \sep Least-square solution \sep \vspace{0.5em}
\noindent\textit{MSC : } 65J10 , 65F05 ,15B33
\end{keywords}

\maketitle

\section{Introduction}\label{sec:introduction}

Obtaining a solution to a linear system of equations is a central task in numerical linear algebra and serves as a foundational component across a broad spectrum of quantitative disciplines. These systems, often expressed in the canonical form $Ax=b$, arise in numerous scientific, engineering, and data-driven applications, including but not limited to computational physics, signal and image processing, optimization, computer vision, machine learning, and numerical simulations. While solving a linear least-squares problem by computing the Moore-Penrose pseudoinverse \cite{svd} has a computational cost that scales cubically, the system in-hand may enjoy a special structure (for example, a system with a partial isometry) that can be utilized to evaluate the least-squares solution at a significantly lower computational cost.

A matrix $A \in \mathbb{C}^{m \times n}$ is said to be a partial isometry if $AA^*A=A$ \cite{partialisometryarxiv}, $A^*$ being the conjugate-transpose of $A$. We present an $O(mn)$ least-squares direct linear solver for the linear system $Ax=b$ where $A$ can be a \emph{scaled} partial isometry. Partial isometries represent a significant and appealing class of operators. An operator is called a partial isometry if it acts as an isometry when restricted to the orthogonal complement of its null space. These operators are fundamental in the field of operator theory; for example, they are essential in the polar decomposition of general operators and serve as a foundation for the theory of von Neumann algebras \cite{partialisometryelsiver}. Common instances of partial isometries include orthogonal projections, unitary operators, isometries, and co-isometries \cite{partialisometryelsiver}. Partial isometries describe certain quantum measurements and state transformations, especially when dealing with non-orthogonal states or incomplete measurements \cite{partialisometryieee}. In signal processing and frame theory, they are used in the construction of tight frames, filter banks, and optimal measurement vectors, where partial isometries ensure energy-preserving transformations on subspaces \cite{partialisometryieee}. In finite dimensions, partial isometries are used to study matrix factorizations, similarity, and unitary equivalence problems \cite{partialisometryarxiv}.

The paper is organized as follows: \cref{subsec:notation} provides the details of the notations used in this paper. A few useful properties of partial isometries are outlined in \cref{sec:properties}. \Cref{sec:methods} contains the proposed algorithm, its applicability, and its theoretical analysis. Experiments were conducted to validate the proposed algorithm, and the corresponding results can be seen in \cref{sec:numerical_experiments}. Finally, we conclude the paper with some remarks in \cref{sec:conclusion}.

\subsection{Notations Used}\label{subsec:notation}
The symbol $A^{\dagger}$ denotes the Moore–Penrose pseudo-inverse \cite{svd} of matrix $A \in \mathbb{C}^{m \times n}$. $A^*$ denotes the conjugate-transpose of $A$. $A^{(i)}, A_{(i)}$ denotes the $i^{\text{th}}$ column and row of the matrix, $A$ respectively. Unless specifically mentioned, $\|.\|$ as well as $\|.\|_2$ denote 2-norm, and $\|.\|_F$ denotes Frobenius norm. For two compatible column vectors $v_1,v_2$, we define the inner product as $\langle v_1, v_2\rangle = v_1^*v_2$.

\section{Properties of Partial Isometries}\label{sec:properties}

\begin{definition}
    A matrix $A \in \mathbb{C}^{m \times n}$ is siad to be a partial isometry if $AA^*A=A$ \cite{partialisometryarxiv}.
\end{definition}

While the above definition is brief, it offers limited insight into the concept of a partial isometry. Nevertheless, it does imply connections to unitary matrices, orthogonal projections, and the Moore–Penrose pseudoinverse—associations that are both meaningful and worth exploring

\begin{theorem}\label{thm:thm1}
    A matrix $A \in \mathbb{C}^{m \times n} $ is a partial isometry iff all its non-zero singular values (if any) are 1.
\end{theorem}
\begin{proof}
    ($\Rightarrow$) To show that if $A$ is a partial-isometry then all its non-zero singular values (if any) are 1.
    
    If $A$ is a zero matrix, then the claim trivially holds. Let's say, rank of $A$ be $r > 1$. Then the partial singular-value decomposition of $A$ can written as $A=U\Sigma V^*$, where $U \in \mathbb{C}^{m \times r}, \Sigma \in \mathbb{R}^{r \times r}, \text{and} \quad V \in \mathbb{C}^{n \times r}$. Now $A$ being a partial isometry: $AA^*A=A$ holds and thus 
    \begin{align*}
        &(U\Sigma V^*)(U\Sigma V^*)^*(U\Sigma V^*)=(U\Sigma V^*)\\
        &\implies U\Sigma^3V^* =U\Sigma V^*\\
        &\implies U^*(U\Sigma^3V^*)V = U^*(U\Sigma V^*)V\\
        &\implies \Sigma^3=\Sigma
    \end{align*}
    Now $\Sigma \in \mathbb{R}^{r \times r}$ being a full-rank diagonal matrix, from above it is clear that the diagonal entries of $\Sigma$ are all 1, as by convention non-zero singular values are positive. And hence all the non-zero singular values of $A$ in this case are 1.

    ($\Leftarrow$) To show that if all the non-zero singular values (if any) of $A$ are 1, then it is a partial isometry.

    If $A$ has no non-zero singular values, then trivially it is a partial isometry. Now, let's say $A$ has $r>1$ non-zero singular values, all of which are 1. Then, using partial singular-value decomposition $A=U\Sigma V^*$, where $\Sigma =I_r$, an $r \times r$ identity matrix. Thus, $A=UV^*$. Now, $AA^*A=(UV^*)(UV^*)^*(UV^*)=UV^*=A$. Thus, A is a partial isometry.
\end{proof}

\begin{theorem}\label{thm:thm2}
    For $A \in \mathbb{C}^{m \times n} $, the following statements has relation $(a) \implies (b)$ and $(a) \implies (c)$.\\
 (a)$A$ is a partial isometry\\
 (b)$A=WP$, where W is a partial isometry and P is an orthogonal projection.\\ 
 (c)$A=QW$, where W is a partial isometry and Q is an orthogonal projection.
\end{theorem}
\begin{proof}
    If $A$ is a zero-matrix, then the proofs for both the relations $(a) \implies (b)$ and $(a) \implies (c)$ are trivial.
    
    To show $(a) \implies (b)$ :\\
     Let $A=U \Sigma V^*$ be a partial singular-value decomposition of the non-zero partial isometry $A$. Then $A=(UV^*)(V\Sigma V^*)$, and it is easy to see that $W=UV^*$ is a partial isometry. Following \cref{thm:thm1} $A$ being a partial isometry $\Sigma$ is an identity matrix and thus $P=V\Sigma V^*=VV^*$ is an orthogonal projection.

    To show $(a) \implies (c)$ :\\
    Similarly, using the partial SVD of $A$, we have $A= U \Sigma V^* = (UU^*)(U\Sigma V^*)$. It is trivial that $Q=UU^*$ is an orthogonal projection and $W=U\Sigma V^*$ is a partial isometry (using the information of $\Sigma$ from \cref{thm:thm1}).
    \end{proof}

\cite{partialisometryarxiv, partialisometryelsiver,partialisometrynote,f1,f3,f5,f6,a1} can be referred for more details about partial isometries.

\section{Solution of a Scaled Partial-Isometric Linear System}\label{sec:methods}
Here we outline the $O(mn)$ procedure to evaluate the minimum 2-norm solution of a system $Ax=b$ where A is a scaled partial isometry.

\begin{theorem}\label{thm:direct_solver}
For a matrix $A \in \mathbb{C}^{m\times n}$ with at least one non-zero singular value and all non-zero singular values being the same, the least square solution of the system $Ax=b$ is given by: 
\begin{equation}
x_i =
\begin{cases}
\frac{(\langle A^{(i)}, b\rangle)^*(\langle A^{(i)}, b\rangle)}{(A^{*}b)^{*}A^{*}A^{(i)}} & \text{if } b^{*}AA^{*}A^{(i)} \neq 0 \\
0 & \text{otherwise}
\end{cases}
\label{eq:a1}
\end{equation}
\end{theorem}
\begin{proof}
The case where A does not have a non-zero singular value is trivial.

Let the partial singular value decomposition of A be $A=U\Sigma V^{*}$, where $U \in \mathbb{C}^{m\times r}, \Sigma \in \mathbb{R}^{r\times r}, V \in \mathbb{C}^{n\times r}$. Let $\alpha$ be the non-zero singular values. Then the least square solution of the system of equations $Ax=b$ is given as $x= V\Sigma ^{-1}U^{*}b$ and thus 
\begin{equation}
x_i= \langle e^{(i)},x\rangle=e^{(i)^*}x=V_{(i)}\Sigma ^{-1}U^{*}b=\frac{1}{\alpha}V_{(i)}U^{*}b
\label{eq:a2}
\end{equation}
[ $e^{(i)}$ being an $n$-dimensional canonical-basis vector with the $i^{th}$ entry as 1 and all other entries as 0.]\\

Now, \begin{align}
\label{eq:a3.1}
& A^{(i)}=U\Sigma V_{(i)}^{*}\\ 
\label{eq:a3.2}
 & A^{(i)^{*}}b=V_{(i)}\Sigma U^{*}b\\
 \label{eq:a3.3}
 & A^{*}b=V\Sigma U^{*}b\\
 \label{eq:a3.4}
 & A^{*}A^{(i)}=V \Sigma ^{2}V_{(i)}^{*}
\end{align}

The condition $b^{*}AA^{*}A^{(i)} = 0$ $ \implies b^{*}U\Sigma V^{*}V\Sigma U^{*}U\Sigma V_{(i)}^{*} =0$, using the relations \eqref{eq:a3.1}- \eqref{eq:a3.4}.

\begin{align*}
 & \implies b^{*}U\Sigma ^{3} V_{(i)}^{*} =0 \\
& \implies \alpha ^3 b^*UV_{(i)}^*=0 \\
& \implies \alpha ^4 \frac{1}{\alpha}V_{(i)}U^{*}b=0\\
& \implies \alpha^4 x_i = 0, \text{ using } \eqref{eq:a2},
\end{align*}
which validates our formulation of considering $x_i=0$ when $b^{*}AA^{*}A^{(i)} = 0$ and $\alpha \neq 0$ in $\eqref{eq:a1}$.

Now, when $b^{*}AA^{*}A^{(i)} \neq 0$, the proposed solution
\begin{align*}
    x_i &= \frac{(\langle A^{(i)}, b\rangle)^*(\langle A^{(i)}, b\rangle)}{(A^{*}b)^{*}A^{*}A^{(i)}} \text{ where using relations (3)-(6) above,}\\
    &=\frac{\alpha^ 2b^*UV_{(i)}^{*}V_{(i)}U^{*}b}{\alpha^3 b^*UV_{(i)}^*}=\frac{1}{\alpha}V_{(i)}U^{*}b
\end{align*}
thus satisfying the required relation \eqref{eq:a2} for the solution.
\end{proof}

\textbf{Note-1 :} The cost of computation of $A^{*}b$ is $O(mn)$, and of $((b^{*}A)A^{*})A$ is also $O(mn)$ using three such matrix-vector products, and thus we can evaluate all components of the least square solution in $O(mn)$ computations using the relation in $\Cref{thm:direct_solver}$. Further, it is to be noted that by replacing $AA^*A^{(i)}$ with $A^{(i)}$ in \eqref{eq:a1} when the exact partial isometry holds, we get the solution $x=A^*b$, the least square solution for a partial isometric system that one can verify as $A(A^*b)=AA^*(Ax)=Ax$ given $AA^*A=A$. Whereas, the proposed solution applies to \emph{scaled} partial isometric systems with an unknown scaling factor as well.

%Partial isometry is invariant under unitarily similar transformation and thus the \cref{thm:direct_solver} is applicable for the systems which are unitarily similar to a sacled partial isometry.

In the following part of the paper, we continue to establish some additional properties of scaled partial isometries. Though the definition in \cref{sec:introduction} is complete and sufficient, these properties may be useful in identifying its application in a practical setting.

\begin{theorem}\label{thm:similarity}
Let \( A \in \mathbb{R}^{n \times n} \). Then, $A$ is similar to  $\mu Q$  for some $\mu> 0$  and an orthogonal  $Q$ iff $A$  is diagonalizable and all eigenvalues of A  have magnitude $\mu$.
\end{theorem} 
% \begin{proof}

% Suppose \( A \) is similar to \( \lambda Q \), where \( \lambda > 0 \) and \( Q \) is an orthogonal matrix. That is, there exists an invertible matrix \( P \) such that
% \[
% A = P (\lambda Q) P^{-1}.
% \]

% Note that since \( Q \) is orthogonal (and a normal matrix), it is diagonalizable with complex eigenvalues \(\mu_1, \dots, \mu_n \) satisfying \( |\mu_i| = 1 \). Therefore, \( \lambda Q \) is diagonalizable as \(S \Lambda S^{-1}\) with eigenvalues \( \lambda \mu_i \), and so \( A \) is diagonalizable as \((PS) \Lambda (PS)^{-1}\) with eigenvalues \( \lambda \mu_1, \dots, \lambda \mu_n \), each with magnitude \( \lambda \) due to the given similarity transformation \( \lambda Q\).

% Now to show that the converse is true as well, let \( A \) be diagonalizable using an invertible matrix $W$ and all its eigenvalues \( \lambda_1, \dots, \lambda_n \) satisfy \( |\lambda_i| = \lambda > 0 \). Then, we need to show that there exists an invertible matrix \( P \) such that
% \[
% A = P \lambda Q P^{-1}.
% \]

% Define
% \[
% \mu_i = \frac{\lambda_i}{\lambda}, \quad \text{so that } |\mu_i| = 1.
% \]

% Then
% \[
% \text{diag}(\lambda_1, \dots, \lambda_n) = \lambda \cdot \text{diag}(\mu_1, \dots, \mu_n),
% \]
% and hence
% \[
% A = W \, \lambda \cdot \text{diag}(\mu_1, \dots, \mu_n) \, W^{-1} = \lambda \cdot \left( PS \, \text{diag}(\mu_1, \dots, \mu_n) \, (PS)^{-1} \right).
% \]

% where we define for some $S$
% \[
% Q := S \, \text{diag}(\mu_1, \dots, \mu_n) \, S^{-1}.
% \]

% Thus \( A = P\lambda Q P^{-1}\). Moreover, as \( A \in \mathbb{R}^{n \times n} \) and the eigenvalues \( \mu_i \) are closed under complex conjugation, \( Q \in \mathbb{R}^{n \times n} \) is orthogonal.
% \end{proof}
\begin{proof}
($\Rightarrow$) Suppose $A$ is similar to $\mu Q$ for some $\mu > 0$ and orthogonal $Q$, i.e., $A = P^{-1} (\mu Q) P$ for some invertible $P$. 

Since $Q$ is orthogonal, it is diagonalizable over $\mathbb{C}$, and so is $\mu Q$. Therefore, $A$ is diagonalizable. The eigenvalues of $Q$ all have modulus $1$, so the eigenvalues of $\mu Q$ (and hence of $A$) all have modulus $\mu$.

($\Leftarrow$) Suppose $A$ is diagonalizable and all its eigenvalues have modulus $\mu > 0$. Then $A$ is similar to a diagonal matrix $D = \mathrm{diag}(\lambda_1, \ldots, \lambda_n)$ with $|\lambda_i| = \mu$ for all $i$. Each $\lambda_i$ can be written as $\mu e^{i\theta_i}$. Over $\mathbb{R}$, complex eigenvalues occur in conjugate pairs, so $D$ is similar over $\mathbb{R}$ to a block-diagonal matrix with real blocks (for real eigenvalues) and $2 \times 2$ rotation blocks (for complex conjugate pairs), which is the canonical form of a real orthogonal matrix. Thus, $A$ is similar to $\mu Q$ for some real orthogonal $Q$.
\end{proof}

Next, we show that for any unitarily diagonalizable matrix whose non-zero eigenvalues are the same up to a multiplication of $\pm 1$, we can have a similar result as in \cref{thm:direct_solver} to evaluate the least-squares solution of the corresponding linear system. These kinds of matrices (and any of their unitary similarity) are also scaled partial isometries, because $A=cUDU^*$, with unitary matrices $U$ and the non-zero diagonal entries of $D \in \{\pm 1\}$. It is easy to see that $P=UDU^*$ is a partial isometry and thus $A=cP$ is a scaled partial isometry.

\begin{theorem}\label{thm:thm3}
    Let $A \in \mathbb{C}^{m \times m}$ be unitarily diagonalizable as $A= UDU^*$, with all the non-zero diagonal entries being the same up to a multiplication of $\pm 1$. Then, $A^\dagger=UD^\dagger U^*$.
\end{theorem}
\begin{proof}
    Let B=$UD^\dagger U^*$. We know $U^*U=I$ and $(DD^\dagger)^*=DD^\dagger$. Hence,
    \begin{align}
        & (AB)^*=(UDU^*UD^\dagger U^*)^*=U(D D^\dagger)^* U^*=UDD^\dagger U^*=AB. \label{eq:thm3_1}\\
        & (BA)^{*}=(UD^\dagger U^*UDU^*)^*=U(D^\dagger D)^*U^*=UD^\dagger DU^*=BA. \label{eq:thm3_2}\\
        & ABA=(UDU^*)(U D^\dagger U^*)(UDU^*)=UDD^\dagger DU^*=UDU^*=A,\label{eq:thm3_3} \text{ using $DD^\dagger D=D$.}\\
        & BAB=(UD^\dagger U^*)(UDU^*)(UD^\dagger U^*)=UD^\dagger DD^\dagger U^*= UD^\dagger U^*=B,\label{eq:thm3_4} \text{ using $D^\dagger D D^\dagger =D^\dagger$.}
    \end{align}
    From \eqref{eq:thm3_1}-\eqref{eq:thm3_4}, we know that $B$ is Moore-Penrose pseudoinverse of $A$. Thus, $A^\dagger = B = U D^\dagger U^*$.
\end{proof}

For such unitarily diagonalizble matrices $A=U_1DU_2^*$ (where$U_1,U_2$ are unitary), as mentioned in \cref{thm:thm3}, we can ensure that all the non-zero diagonal entries of $D$ have the real part as non-negative. For the instances, $Re(D_{i,i})<0$, one can flip the sign of $i^{th}$ column of $U_1$, and also of $i^{th}$ row of $D$ to obtain $\tilde{U}, \tilde{D}$ respectively, such that $U_1D=\tilde{U}\tilde{D}$. Now we can use \cref{thm:direct_solver} to evaluate the least-square solution $x_\star=U\tilde{D}^\dagger\tilde{U}^*b$ in $O(mn)$ computations for any such scaled systems (and their unitarily similar systems) in \cref{thm:thm3}.

\textbf{Note-2 :} For a linear system of equations represented as \( Ax = b \), if \( A \) is a block diagonal matrix and each block is a scaled partial isometry, we can still apply Theorem \ref{thm:direct_solver} to each block to obtain the minimum 2-norm solution. It is not necessary for all blocks to be scaled by the same scalar \( c \). Consequently, these types of matrices \( A \) may have different non-zero singular values and eigenvalues, and they are not necessarily the same up to a multiplication of \( \pm 1 \).

\section{Numerical Experiments }\label{sec:numerical_experiments}
 
For \Cref{tab:matrix_with_equal_singular_values}, we generated real matrices using \Cref{algo:matrix_with_equal_singular_values} with entries of $X$, $Y$ from a standard normal distribution. We generated $t \in \mathbb{R}^{n \times 1}$ with entries from a standard normal distribution and $b=At$. For \Cref{tab:matrix_with_equal_singular_values1}, we generated complex matrices by having $X=X_R + iX_I, Y=Y_R + iY_I$ in \Cref{algo:matrix_with_equal_singular_values}. We also generated $t \in \mathbb{C}^{n \times 1}$ through $t_R + it_I$ where we evaluate $b=At$. All random variables $X_R$, $X_I$, $Y_R$, $Y_I$, $t_R$, $t_I$ were independently drawn from a standard normal distribution. In both cases, the true solution $x_\star$ was obtained as $A^\dagger b$. For every fixed size, we averaged over 150 trials and reported the mean value of $\| x-x_\star \|_2$. All experiments have been executed in MATLAB 2024b on a computer with Intel(R) Core(TM) i9-14900K @ 6.0 GHz, 62.00 GB RAM, and 16.00 GB memory.

\Cref{tab:matrix_with_equal_singular_values}and \Cref{tab:matrix_with_equal_singular_values1} presents these experimental results demonstrating the validity of \Cref{thm:direct_solver}. We also know that the algorithm needs only three matrix-vector products in the evaluation of the true solution, which are numerically stable operations, and hence the method is numerically stable.

 \begin{table}%[]
\caption{Validation of \Cref{thm:direct_solver} with \emph{real} matrices generated using  \Cref{algo:matrix_with_equal_singular_values}}\label{tab:matrix_with_equal_singular_values}
\begin{tabular*}{\tblwidth}{@{}LL@{}LL@{}LL@{}LL@{}}
\toprule
  s & 10 & 10 & 10 & 10 \\
        
  Size & 10000 $\times$ 10000 & 10000 $\times$ 2000  & 5000 $\times$ 10000 & 5000 $\times$ 5000\\
         
  r & 2000 & 400 & 2000 & 5000\\
\midrule
 $\|x-x_\star\|_2$  & 1.81e-13 & 3.34e-14  & 1.79e-13 & 1.97e-13\\
\bottomrule
\end{tabular*}
\end{table}

\begin{table}
\caption{Validation of \Cref{thm:direct_solver} with \emph{complex} matrices generated using  \Cref{algo:matrix_with_equal_singular_values}}\label{tab:matrix_with_equal_singular_values1}
\begin{tabular*}{\tblwidth}{@{}LL@{}LL@{}LL@{}LL@{}}
\toprule
  s & 10 & 10 & 10 & 10 \\
        
  Size & 10000 $\times$ 10000 & 10000 $\times$ 2000  & 5000 $\times$ 10000 & 5000 $\times$ 5000\\
         
  r & 2000 & 400 & 2000 & 5000\\
\midrule
 $\|x-x_\star\|_2$  &2.25e-13  &1.20e-13   &6.79e-15  &5.23e-14 \\
\bottomrule
\end{tabular*}
\end{table}

\begin{algorithm}
\caption{Generate Random Matrix with Equal Non-zero Singular Values}\label{algo:matrix_with_equal_singular_values}
\begin{algorithmic}[1]
\Require Integers $m, n, r$ with $r \leq \min(m,n)$, scalar $s$
\Ensure Matrix $A \in \mathbb{C}^{m \times n}$ with rank $r$ and all non-zero singular values equal to $s$
\State Generate a random matrix $X \in \mathbb{C}^{m \times m}$
\State Compute QR decomposition: $X = QR$, set $U \gets Q$
\State Generate a random matrix $Y \in \mathbb{C}^{n \times n}$
\State Compute QR decomposition: $Y = QR$, set $V \gets Q$
\State Initialize $\Sigma \in \mathbb{R}^{m \times n}$ as a zero matrix
\For{$i = 1$ to $r$}
    \State $\Sigma_{i,i} \gets s$
\EndFor
\State Compute $A \gets U \cdot \Sigma \cdot V^*$
\State \Return $A$
\end{algorithmic}
\end{algorithm}

\section{Conclusion}\label{sec:conclusion}
We presented a $O(mn)$ direct least-squares solver for the system of linear equations $Ax=b$ where A is a scaled partial isometry, and also where A is a block diagonal matrix with each block being a scaled partial isometry. Numerical experiments demonstrate the validity of the proposed method.  

%\clearpage %%Remove this from your manuscript

% To print the credit authorship contribution details
%\printcredits

%% Loading bibliography style file
%\bibliographystyle{model1-num-names}
\bibliographystyle{cas-model2-names}

% Loading bibliography database
\bibliography{reference}

% Biography
%\bio{}
% Here goes the biography details.
%\endbio

%\bio{pic1}
% Here goes the biography details.
%\endbio

\end{document}

